\documentclass[instructions.tex]{subfiles}
\begin{document}
\chapter{Le respect humain, obstacle à la conversion}

\primo Le respect humain est un obstacle à la conversion, et souvent une source de réprobation.

Voici la pierre de scandale pour plusieurs. Si je change, si je ne vois plus cette personne, si je ne vais plus dans cette compagnie, si je quitte cette maison, si je ne tire pas vengeance de cette injure, si je ne fais pas comme les autres, que dira-t-on de moi ?

Que dira-t-on de moi ? Vous craignez donc plus les discours du monde que les jugemens de Dieu, et par là vous faites injure au Très-Haut. Lorsque Lucifer voulut entraîner les bons anges dans sa révolte, ils s’écrièrent : Qui est-ce donc qui est comparable à\footnote{Quærens me sedisti lassus, redemisti crucem passus, tantus labor non sit cassus. Prose.} Dieu ? Servez-vous de cette même pensée contre le respect humain.

Que dira-t-on de moi ? Regardez-vous donc comme un opprobre de servir Dieu\footnote{Quis ut Deus ? Humbert. Pensées.} ? On se fait honneur de servir un prince, et vous regardez comme un déshonneur de servir Jésus-Christ. Le plus vil artisan se fait honneur d’exercer sa profession, quelque misérable qu’il soit ; et vous, chrétiens, vous rougissez du christianisme, vous n’osez paraître disciples de Jésus-Christ.

Que dira-t-on de moi ? Mais qu’importe quoi qu’on en dise, pourvu que vous fassiez votre devoir ? Pourquoi vous mettre en peine du monde ? Vous ne gagnez rien à ses applaudissemens ; vous ne perdez rien à sa censure. Êtes-vous obligé de rendre compte aux hommes de vos actions ? sont-ils vos juges ? en attendez-vous votre récompense ? Dieu seul doit vous juger : pourquoi vous embarrasser du jugement des créatures ?

Que dira-t-on de moi ? Et qu’en dit-on à présent ? Quelques mondains vous applaudissent sur vos désordres, tandis que dans leur cœur ils vous méprisent. Mais les gens de bien que pensent-ils de vous ? Que dit-on dans le public de vos fréquentations, de votre vanité, de vos intrigues, de vos débauches, de vos scandales ? Voilà les discours que vous devez craindre.

Vous êtes si aveuglé, que vous n’avez point de honte de paraître sans pudeur et sans religion, tandis que vous avez honte de paraître vertueux. Si vous devez avoir de la honte, c’est de vous-même. Magdeleine, ouvrant les yeux sur les misères de son âme, fut chargée d’une telle confusion, que, n’osant paraître devant le Sauveur, elle se prosterna à ses pieds, en versant des larmes de repentir et de douleur. Quoiqu’elle fût une fille de qualité, elle ne rougit point de paraître pénitente, de renoncer à ses intrigues devant une assemblée de personnes distinguées. Telles doivent être les dispositions d’une âme pour se convertir.

\secundo Le respect humain, en mettant obstacle à la conversion, est pour plusieurs une source de réprobation. Celui, dit Jésus-Christ, qui rougit de m’appartenir et de pratiquer ma doctrine, le Fils de l’Homme rougira de le reconnoître au nombre de ses élus\footnote{Qui me erubuerit et sermones meos, hunc Filius hominis erubescet. Luc. 9.}. C’est, au contraire, une marque de prédestination, de pratiquer avec courage les maximes de Jésus-Christ : Celui qui me connoîtra devant les hommes, dit le Sauveur, je le reconnoîtrai devant mon Père\footnote{Omnis ergo qui confitebitur me coram hominibus, confitebor et ego eum coram Patre meo qui in cœlis est. Matth. 10. 32.}.

C’est avoir l’âme basse, dit saint Martin de Brague, de n’oser être sage, parce que les fous s’en moquent. Leurs discours passeront, mais les jugemens de Dieu ne passeront pas. Ces mondains vous tireront-ils de l’enfer, où votre complaisance et votre lâcheté vous auront précipité ? Ne craignez point de passer pour singulier aux yeux des hommes ; mais craignez d’être réprouvé de Dieu. Pourvu que Dieu soit content de vous, soyez peu en peine des railleries du monde. Il est glorieux, dit ce saint évêque, d’être méprisé des méchans\footnote{Malis displicere est laudari.}.

Si votre conduite leur paraît ridicule et bizarre, qu’importe, pourvu qu’elle serve à vous sauver ?

Quoique vous deviez mépriser les discours du monde quand il s’agit de votre devoir, vous ne devez pas donner imprudemment occasion à la satire et aux plaintes d’autrui, ni refuser à vos amis, à vos parens, à vos supérieurs, ce que la complaisance ou le devoir exige dans les choses permises ; mais si ce qu’ils demandent est contre la loi de Dieu, il n’y a personne au monde pour qui vous deviez avoir de la complaisance au préjudice de votre conscience. Quelque autorité que l’on ait sur vous, quelque liaison que vous ayez de parenté, d’amitié ou d’intérêt, il vaut mieux déplaire cent fois à tous les hommes que de déplaire une seule fois à Dieu.

O respect humain ! que tu as perdu d’âmes ! Dieu veuille que le sanctuaire en soit exempt ! Quel malheur pour les ministres sacrés, par des vues humaines, d’abandonner la discipline, et de flatter les passions ! Si je craignais les hommes, dit saint Paul, je ne serais pas serviteur de Jésus-Christ\footnote{Si adhuc hominibus placerem, Christi servus non essem. Gal. 1. 10.}.

\section{Résolutions}

\primo Je me mettrai désormais toute ma gloire à servir Dieu.

\secundo Je ne craindrai point les discours du monde, lorsqu’il s’agira de mon salut.

\tertio Je me conduirai avec prudence, pour ne point donner sujet aux médisances, sans jamais sacrifier ma conscience au respect humain.

Elevez en mon coeur, Ô mon Dieu! au-dessus de toute crainte humaine, afin que je demeure constamment attaché à votre service, et que je ne trahisse jamais plus les intérêts de votre gloire par respect humain 

\end{document}
